% !TEX root = ../eigenfaces_algorithm.tex
%=====================================================
\section{数据组织：矩阵、均值脸与中心化}

\subsection{训练集与划分}

把 400 张照片按人分层划分：每人 10 张里固定取 7 张进训练集、3 张进测试集。
训练集共 $n=280$ 张，堆成一个数据矩阵：
\[
\mat{A}=\begin{pmatrix}
\rule{4pt}{9pt}\ \mathbf{a}_1^{\mathsf{T}}\ \rule{4pt}{9pt}\\
\rule{4pt}{9pt}\ \mathbf{a}_2^{\mathsf{T}}\ \rule{4pt}{9pt}\\
\vdots\\
\rule{4pt}{9pt}\ \mathbf{a}_{280}^{\mathsf{T}}\ \rule{4pt}{9pt}
\end{pmatrix}\in\R^{280\times 4096},
\qquad n=280,\quad d=4096 .
\]
每一\textbf{行}是一张脸（拉直后的向量转置），每一\textbf{列}是“某个像素位置在 280 张脸里的取值”。

\subsection{均值脸与中心化}

对 280 张训练脸逐像素取平均，得到\textbf{均值脸} $\bar{\mathbf{a}}=\frac1n\sum_{j}\mathbf{a}_j$。
把每张脸减去它：
\[
\mat{X}=\mat{A}-\mathbf{1}\,\bar{\mathbf{a}}\mT ,
\]
即 $\mat{X}$ 的每一行都是“这张脸减去平均长相”的\textbf{残差脸}。
（$\mathbf{1}$ 是 280 维全 1 向量，作用是把同一个均值脸复制 280 份逐行减掉。）

\begin{why}
如果不减均值会怎样？所有像素都非负、整体亮度贡献了一个巨大的“公共成分”，
数据云的重心远离原点。此时方差最大的方向几乎就是“平均亮度方向”，
它对区分谁是谁毫无帮助。中心化把云的重心挪到原点，
剩下的方差才反映“人与人、脸与脸的差异”。
\end{why}

\begin{why}
均值只用训练集算。测试集若用自己的均值，等于让模型提前“看了”测试数据（信息泄漏），
而且测试脸的坐标系会和训练出的特征脸对不上。这是本题第一坑。
\end{why}

\subsection{一个决定一切的不等式}

中心化后每行和为定值（都减了同一个均值），280 行中只有 $n-1=279$ 行独立，故
\[
\operatorname{rank}(\mat{X})\le n-1=279\ll 4096 .
\]
\begin{intuition}
280 张残差脸不管看起来多不一样，它们全体都躺在 4096 维空间里一个
\textbf{至多 279 维的薄板}上。也就是说，4096 个维度里至少有 3817 个是“空房间”。
我们不需要在 4096 维里找东西，只需要描述这块薄板——这就是后面一切“小矩阵技巧”的底气。
\end{intuition}
